\documentclass[12pt]{article}
\usepackage[utf8]{inputenc}
\usepackage[T1]{fontenc}
\usepackage[margin=1in]{geometry}
\usepackage{setspace}
\usepackage{xcolor}
\usepackage[colorlinks=true, urlcolor=blue!60!black, linkcolor=blue!60!black]{hyperref}
\setlength{\parindent}{2em}
\setlength{\parskip}{1.25em}
\renewcommand{\baselinestretch}{1.0}
\begin{document}
\begin{center}
{\LARGE \textit{Shurtleff v. City of Boston}, 596 U.S. 243 (2022)\par}
\vspace{0.8em}
{\large Case Brief\par}
\end{center}
\vspace{1em}

\section{Memory Jogger}
Boston had approved 284 straight private requests to fly flags on a City Hall flagpole but refused Camp Constitution's Christian flag; the Court held 9--0 that, because Boston did not actively control the flags' messages, the program was private speech, and the refusal was unconstitutional viewpoint discrimination (\href{https://supreme.justia.com/cases/federal/us/596/20-1800/}{Justia}).

\section{Detailed Case Facts}
Outside Boston City Hall on City Hall Plaza---a 7-acre public gathering space that Boston acknowledged was a public forum---stand three 83-foot flagpoles. The first pole flies the United States flag (with the POW/MIA banner), the second flies the Massachusetts flag, and the third ordinarily flies the City of Boston's flag. Since at least 2005, however, Boston allowed private groups to hold ceremonies on the plaza and temporarily hoist a flag of their choosing on the third flagpole in place of the city flag. Between 2005 and 2017, Boston approved 284 consecutive flag-raising ceremonies involving roughly 50 unique flags---mostly foreign national flags, alongside civic and private banners such as the Pride Flag, an emergency medical services banner, and a local community bank's flag (Metro Credit Union). Boston had no written policy governing flag content and had never rejected a single application.

In July 2017, Harold Shurtleff, director of Camp Constitution, applied to hold a September ceremony on City Hall Plaza to ``commemorate the civic and social contributions of the Christian community'' and to raise what he described as the ``Christian flag'' (a red cross on a blue field against a white background). Gregory T. Rooney, commissioner of Boston's Property Management Department, denied the request---the first denial in the program's history---solely because he feared flying a religious flag at City Hall would violate the Establishment Clause. Rooney told Shurtleff the group could hold its ceremony only if it flew a different flag.

\section{Procedural History}
Shurtleff and Camp Constitution sued Boston and Rooney in the U.S. District Court for the District of Massachusetts, alleging violations of the First Amendment's Free Speech Clause. After denying preliminary injunctive relief (337 F. Supp. 3d 66 (D. Mass. 2018), aff'd, 928 F.3d 166 (1st Cir. 2019)), the District Court granted summary judgment for Boston on a stipulated record, holding that the third-pole displays constituted government speech. The First Circuit affirmed (986 F.3d 78 (1st Cir. 2021)). The Supreme Court granted certiorari and unanimously reversed.

\section{Judicial Votes}
The Supreme Court unanimously \textbf{reversed and remanded} (9--0). The case was argued on January 18, 2022, and decided on May 2, 2022 (\href{https://supreme.justia.com/cases/federal/us/596/20-1800/}{Justia}; \href{https://www.law.cornell.edu/supremecourt/text/20-1800}{Cornell LII}):

\begin{itemize}
    \item \textbf{Majority Opinion:} Justice Breyer delivered the opinion of the Court, joined by Chief Justice Roberts and Justices Sotomayor, Kagan, Kavanaugh, and Barrett.
    \item \textbf{Concurring Opinions:} Justice Kavanaugh filed a concurring opinion. Justice Alito filed an opinion concurring in the judgment, joined by Justices Thomas and Gorsuch. Justice Gorsuch filed an opinion concurring in the judgment, joined by Justice Thomas.
    \item \textbf{Dissenting Opinions:} None.
\end{itemize}

\section{Holding}
\textbf{Questions Presented:}

\begin{enumerate}
  \item Did Boston's program allowing private groups to raise flags of their choosing on a City Hall Plaza flagpole constitute government speech exempt from Free Speech Clause scrutiny, or private speech in a public forum?
  \item If the flag raisings were private speech, did Boston's refusal to fly Camp Constitution's Christian flag based on Establishment Clause concerns violate the Free Speech Clause as impermissible viewpoint discrimination?
\end{enumerate}

\textbf{Holding:} Private speech, and yes (9--0). Applying a holistic, context-driven inquiry weighing history, public perception, and the extent of active government control, the Court held that Boston's flag-raising program did not express government speech. Because Boston opened the flagpole as a forum for private expression, refusing to fly Camp Constitution's flag based on its religious viewpoint violated the First Amendment's Free Speech Clause.

\section{Analysis of Opinions}
\subsection{Majority Opinion (Justice Breyer)}

Writing for a six-Justice majority, Justice Breyer began from first principles: when the government speaks for itself, the Free Speech Clause does not require viewpoint neutrality---democratic accountability at the ballot box checks government messaging (\textit{Pleasant Grove City v. Summum}, 555 U.S. 460 (2009); \textit{Walker v. Texas Div., Sons of Confederate Veterans, Inc.}, 576 U.S. 200 (2015)). When the government invites private participation, however, courts must conduct a ``holistic inquiry''---not a mechanical checklist---examining (1) the history of the expression, (2) the public's likely perception of who is speaking, and (3) the extent to which the government actively shaped or controlled the expression:
\begin{itemize}
  \item \textbf{History of Flag Flying (Favors Boston):} General history supported the City. Flags at seats of government---from national standards and half-staff tributes to foreign flags outside Blair House---have long conveyed official governmental messages, and Boston's third pole usually communicated city messages. Yet general history was only the starting point.
  \item \textbf{Public Perception (Inconclusive):} While a passerby seeing three matching flagpoles outside City Hall might ordinarily assume all three spoke for Boston, the City regularly lowered its own flag for private ceremonies held at the pole's base. A pedestrian watching private citizens conduct an event without city officials might attribute the banner to the group, not the City.
  \item \textbf{Active Government Control (Decisive Against Boston):} The ``most salient feature'' was that Boston ``neither actively controlled these flag raisings nor shaped the messages the flags sent.'' Although Boston scheduled dates and lent a hand crank, it told the public it sought ``to accommodate all applicants'' at its ``public forums,'' asked only for basic contact and event details, never inspected flags beforehand, maintained no written content guidelines, and approved 284 straight requests (including a local credit union's flag). Unlike Pleasant Grove's selective acceptance of permanent park monuments (\textit{Summum}) or Texas's active review and rejection of specialty license plates (\textit{Walker}), Boston's ``come-one-come-all attitude'' mirrored the Patent and Trademark Office's registration of private trademarks in \textit{Matal v. Tam}, 582 U.S. 218 (2017). Breyer noted, however, that nothing barred Boston from changing its policy going forward (citing San Jose's written closed-list flag policy).
  \item \textbf{Viewpoint Discrimination (Part III):} Once Boston's government-speech premise collapsed, its Establishment Clause defense collapsed with it. Because Boston admitted denying Shurtleff's flag solely because it ``promot[ed] a specific religion,'' the denial constituted unconstitutional viewpoint discrimination under \textit{Good News Club v. Milford Central School}, 533 U.S. 98 (2001), and \textit{Rosenberger v. Rector and Visitors of Univ. of Va.}, 515 U.S. 819 (1995).
\end{itemize}

\subsection{Concurring Opinion (Justice Kavanaugh)}

Justice Kavanaugh emphasized in a brief concurrence that the dispute stemmed solely from a city official's mistaken belief that treating religious speech equally in a public program violates the Establishment Clause. Under the Constitution, excluding religious speakers from neutral public forums violates the First Amendment because government may never treat religion as ``second-class.''

\subsection{Concurring in the Judgment (Justice Alito)}

Justice Alito rejected the \textit{Summum}/\textit{Walker} three-factor triad as an unprincipled test susceptible to masking censorship---pointing out that heavy ``control'' is also the hallmark of a royal censor (like Britain's Lord Chamberlain) and that ``public perception'' lets government exploit passerby misattribution. Alito proposed a two-part bright-line rule: expression is government speech \emph{only} if (1) the government purposefully communicates a governmentally determined message through authorized agents (or voluntarily alienated/adopted media), and (2) does so without abridging private speech. Because Boston opened a forum and accepted 284 unvetted private flags (including both the Pride Flag and Ethiopia's flag), it was not speaking for itself.

\subsection{Concurring in the Judgment (Justice Gorsuch)}

Justice Gorsuch traced Boston's error directly to \textit{Lemon v. Kurtzman}, 403 U.S. 602 (1971), and its ``reasonable observer'' endorsement offshoot, which Gorsuch mocked as a ``children's game'' that trapped risk-averse localities in a ``\textit{Lemon} trade''---censoring religious speech to avoid imaginary Establishment Clause suits. Gorsuch urged courts and localities to bury \textit{Lemon} and rely instead on Founding-era ``historical practices and understandings,'' identifying six historical hallmarks of religious establishment (mostly involving coercion or government control over churches) and noting that neutral public-forum access implicates none of them.

\section{Examples: Future Applications}
\begin{enumerate}

\item \textbf{Same-Side (Private Speech; Exclusion Invalid): The Library's Community Heritage Wall.} A municipal library installs a ``Community Heritage Wall'' in its main lobby and adopts a written ordinance stating that the wall is ``not a public forum'' and displays ``official city-endorsed messages celebrating local history.'' In practice, however, the library posts an online form inviting residents to submit plaques honoring local figures, approves 140 consecutive submissions without committee review---including plaques for a local brewery founder and a high-school mascot---and then rejects a mosque's plaque honoring a local Muslim philanthropist on the ground that it might offend secular patrons or violate the Establishment Clause. Under \textit{Shurtleff}, a self-serving written disclaimer does not settle the inquiry when actual practice reveals zero meaningful governmental control over content. Like Boston's 284 unvetted flag raisings, the library's come-one-come-all administration establishes a forum for private speech; excluding the mosque's plaque for its religious viewpoint violates the Free Speech Clause (\textit{Good News Club}), and under \textit{Kennedy v. Bremerton} and Gorsuch's \textit{Shurtleff} concurrence, neutral inclusion poses no Establishment Clause violation.

\item \textbf{Same-Side (Private Speech; Exclusion Invalid): Holiday Displays in the Town Green.} For fifteen years a town has issued a permit to any resident group that asks to place a seasonal display on its town green in December, without reviewing designs; it has approved nativity scenes, menorahs, and a local scout troop's snowman. When an atheist association applies to erect a sign celebrating the winter solstice and questioning religious belief, the town refuses, saying the sign is offensive. As in \textit{Shurtleff}, history alone cannot carry the day when the town has exercised no meaningful control over the displays' messages, so the green hosts private speech. Rejecting one group's display because of its views on religion is viewpoint discrimination that the Free Speech Clause forbids (\href{https://supreme.justia.com/cases/federal/us/596/20-1800/}{Justia}).

\item \textbf{Opposite-Side (Government Speech): A Closed-List Flag Policy.} After \textit{Shurtleff}, a city adopts a written policy that its third flagpole will fly only the city flag or flags chosen by city council resolution to express the city's own messages, and it stops accepting applications from private groups. A religious group later asks to raise its flag and is refused. The policy is the kind of change Justice Breyer said nothing prevented Boston from making: the city now actively selects and controls every flag, the public has no access to the pole, and the history of flags at seats of government already favors the city. The flags are government speech, so the Free Speech Clause does not require the city to accept the group's flag.

\item \textbf{Opposite-Side (Government Speech): Actively Reviewed Specialty License Plates.} A state motor-vehicle board reviews every proposed specialty license-plate design, routinely rejects designs it finds inconsistent with state messaging, and owns the designs it approves. The board rejects a proposed plate bearing an advocacy group's political slogan. Under \textit{Walker v.\ Texas Div., Sons of Confederate Veterans, Inc.}, 576 U.S.\ 200 (2015), which \textit{Shurtleff} distinguished, plates are traditionally used to convey state messages, observers associate them with the State, and the board actively controls their content. All three factors point to government speech, so the board may reject the design without satisfying viewpoint neutrality.

\item \textbf{Fence-Sitter (Close Call): A Juried Public Mural Program.} A city invites artists to propose murals for the walls of city-owned buildings. A city arts commission reviews each proposal for artistic quality and site compatibility and has approved about 90 percent of them, but it has no written guidelines on content. It rejects a proposal featuring religious imagery. The commission's review is more active than Boston's rubber stamp, and murals on public buildings may be read as the city's own message, which points toward government speech under \textit{Summum}. Yet the high approval rate, the absence of content standards, and the program's open invitation resemble \textit{Shurtleff}'s come-one-come-all forum. Because \textit{Shurtleff}'s holistic inquiry gives no fixed weight to each factor, the result could go either way.

\end{enumerate}

\section{Critique}
\begin{itemize}
  \item \textit{The ``Leveling Down'' Paradox:} By confirming that Boston could avoid viewpoint-neutrality obligations simply by adopting a written closed-list policy, \textit{Shurtleff} incentivizes municipalities to shut down public forums altogether---reducing total private speech on public property to avoid hosting controversial flags (such as the Satanic Temple's post-\textit{Shurtleff} application in Boston).
  \item \textit{Control as Both Sword and Shield:} Alito's concurrence exposes a deep tension in Breyer's third factor: the more aggressively a city polices and rejects private speech before it is displayed, the more likely a court is to reward that censorship by calling the program ``government speech.''
\end{itemize}

\section{Key Quotations}
\begin{quote}
``In answering these questions, we conduct a holistic inquiry designed to determine whether the government intends to speak for itself or to regulate private expression. Our review is not mechanical; it is driven by a case's context rather than the rote application of rigid factors.''
\end{quote}
\noindent --- Justice Breyer, 596 U.S. at 252.

\begin{quote}
``All told, while the historical practice of flag flying at government buildings favors Boston, the city's lack of meaningful involvement in the selection of flags or the crafting of their messages leads us to classify the flag raisings as private, not government, speech---though nothing prevents Boston from changing its policies going forward.''
\end{quote}
\noindent --- Justice Breyer, 596 U.S. at 258.

\begin{quote}
``Under the Constitution, a government may not treat religious persons, religious organizations, or religious speech as second-class.''
\end{quote}
\noindent --- Justice Kavanaugh, concurring, 596 U.S. at 261.

\begin{quote}
``To avoid a spurious First Amendment problem, Boston wound up inviting a real one. Call it a \textit{Lemon} trade\ldots This Court long ago interred \textit{Lemon}, and it is past time for local officials and lower courts to let it lie.''
\end{quote}
\noindent --- Justice Gorsuch, concurring in the judgment, 596 U.S. at 280, 288.


\end{document}
